\section{El crecimiento del horizonte temporal y la brecha de confiabilidad}

\subsection{De segundos a casi una hora}

\videofigure{figures/fifty-percent-horizon.jpg}{50\% horizon: crecimiento desde el nivel de segundos de GPT-2 hasta casi una hora en Claude 3.7 Sonnet; la tendencia histórica se aproxima a un crecimiento exponencial.}{00:14:02--00:16:18}

El ajuste presentado en el curso muestra que, en los últimos seis años aproximadamente, el doubling time del time horizon ha sido de unos siete meses. Esta tendencia implica que el agent no solo es más preciso en tareas cortas del mismo tipo, sino que también puede mantener cadenas más largas de planificación y llamadas a herramientas.

Si el horizon crece exponencialmente con el año $y$, puede escribirse como:

$$
H(y)=H_0\,2^{(y-y_0)/\tau},
$$

\begin{itemize}
    \item $H_0$: el time horizon en el año de referencia;
    \item $y-y_0$: los años transcurridos;
    \item $\tau$: el doubling time, que el curso estima en unos siete meses.
\end{itemize}

\begin{warningbox}{No tomes la extrapolación de un historial corto como una predicción certera}
La frecuencia de publicación de los modelos, la composición del conjunto de tareas, el entorno de herramientas y la filtración de datos pueden alterar la curva. El ajuste exponencial describe las muestras pasadas; no garantiza que se mantenga la misma pendiente durante los próximos seis años.
\end{warningbox}

\subsection{50\% y 80\% horizon}

\videofigure{figures/eighty-percent-horizon.jpg}{El 80\% success horizon también mejora entre generaciones, pero su duración absoluta es considerablemente más corta que la del 50\% horizon.}{00:19:42--00:21:42}

\videofigure{figures/horizon-comparison.jpg}{El 50\% horizon de Claude 3.7 es de unos 59 minutos, mientras que el 80\% horizon es de unos 15 minutos, lo que revela una brecha de confiabilidad en tareas de dificultad media.}{00:21:42--00:23:18}

La métrica del 50\% responde a «qué tan largo es lo que hay la mitad de probabilidad de lograr»; la métrica del 80\% se acerca más a lo que exige un despliegue real. Las pendientes de crecimiento de ambas pueden ser similares, pero $H_{80}$ es varias veces más corta, lo que indica que el modelo sigue presentando fallas aleatorias inaceptables con frecuencia en tareas más largas.

\begin{importantbox}{Lo que importa para el despliegue es la reliability frontier}
Un sistema real debería reportar una curva $H_p$ completa, no solo $H_{50}$. Las tareas médicas, financieras o de infraestructura pueden requerir $p\ge 0.99$; en cambio, la generación de borradores de bajo riesgo puede aceptar fallas seguidas de un reintento.
\end{importantbox}

\subsection{Resumen de la sección}

El alcance de tareas autónomas del agent crece rápidamente, pero sigue existiendo una brecha enorme entre «lograrlo alguna vez» y «lograrlo de manera constante». Cuanto más alto es el umbral de confiabilidad exigido, más corta es la duración de las tareas que pueden completarse de forma automática.

